\subsection{Carmichael function}

The Carmichael function $\lambda(n)$ is the smallest positive integer $m$ such that

\begin{equation}
    a^{m} \equiv 1 \pmod{n}
\end{equation}

holds for every integer $a$ coprime to $n$. In other words, it is the exponent of the
multiplicative group of integers modulo $n$.

\subsubsection{Computation from the prime factorization}

For a prime power $p^{r}$,

\begin{equation}
    \lambda(p^{r}) =
    \begin{cases}
        \tfrac{1}{2}\varphi(p^{r}) & \text{if } p = 2 \text{ and } r \geq 3, \\
        \varphi(p^{r})             & \text{otherwise,}
    \end{cases}
\end{equation}

where $\varphi(p^{r}) = p^{r-1}(p-1)$. For a general $n = p_1^{r_1} p_2^{r_2} \cdots p_k^{r_k}$,

\begin{equation}
    \lambda(n) = \operatorname{lcm}\left( \lambda(p_1^{r_1}),\, \lambda(p_2^{r_2}),\, \ldots,\, \lambda(p_k^{r_k}) \right).
\end{equation}

\subsubsection{Relation to Euler's totient}

Since $\lambda(n)$ is the exponent of the multiplicative group modulo $n$ while
$\varphi(n)$ is its order,

\begin{equation}
    \lambda(n) \mid \varphi(n).
\end{equation}

\subsubsection{Example}

For $n = 8$ we have $\varphi(8) = 4$ but $\lambda(8) = 2$ (the special case for powers of
$2$). Indeed $1^{2} \equiv 3^{2} \equiv 5^{2} \equiv 7^{2} \equiv 1 \pmod{8}$.

\subsubsection{First values}

\begin{center}
\begin{tabular}{c|cccccccccccccccc}
$n$           & 1 & 2 & 3 & 4 & 5 & 6 & 7 & 8 & 9 & 10 & 11 & 12 & 13 & 14 & 15 & 16 \\
\hline
$\lambda(n)$  & 1 & 1 & 2 & 2 & 4 & 2 & 6 & 2 & 6 & 4  & 10 & 2  & 12 & 6  & 4  & 4
\end{tabular}
\end{center}
